\section{Introduction}
\label{sec:introduction}

%===========第一段的思考
%一句话介绍背景：LLM-as-a-judge 方法使用广泛，为了让模型按评价标准进行判断，并方便最终分数
%可解释和可审查，很多研究会让其先生成 rationale 再给出分数。而科学写作评价领域，这个更依赖领域知识，且评价包含多个维度，
%更需要引入 rationale 来辅助判断，

% LLM-as-a-judge 已成为开放式生成评价中的常用范式；
% 为促使模型按评价标准进行判断，并提高最终分数的可解释性与可审查性，
% 许多方法要求模型先生成 rationale，再给出分数
% ~\citep{liu-etal-2023-g,kim-etal-2024-prometheus-iclr}。
% 在科学写作评价中，这一设计更具必要性，因为判断通常依赖领域知识、文本证据和多个评价维度；
% ~\citep{sadallah-etal-2025-good,sahinuc-etal-2026-reward,zhang-etal-2026-surveyreview}。

LLM-as-a-judge has become a widely used paradigm for evaluating open-ended generations;
to support criterion-grounded judgments and make final scores more interpretable and auditable,
many methods elicit rationales before scoring
~\citep{liu-etal-2023-g,kim-etal-2024-prometheus-iclr}.
In scientific-writing evaluation, this rationale-first design is particularly relevant
because judgments often depend on domain knowledge, textual evidence and multiple quality dimensions
~\citep{sadallah-etal-2025-good,sahinuc-etal-2026-reward,zhang-etal-2026-surveyreview}.

%=============第二段的思考
%上面介绍了需要引入 rationale 来评价，这里就需要介绍有哪些方法引入 rationale
%实现了好的评测效果。并指出这些方法的问题是太多东西改变了，实验不能完全证明是引入了 rationale 就达到了他们说的效果。
%然后点出两篇经典的文章，具体指出并分析他们的问题是什么。
% sahinuc-etal-2026-reward 文章是两阶段强化学习，让模型自己生成并反思纠正，他的效果是：；
% geval 是通过第三方大模型api先提示词生成 rationale 再加入到最终输入中，让模型进行ing测，他的效果是：

% 现有评价方法已通过提示生成、监督蒸馏和强化学习等不同方式引入 rationale，并普遍报告了更高的人类相关性、judge agreement 
% 或泛化能力~\citep{liu-etal-2023-g,kim-etal-2024-prometheus-iclr,sahinuc-etal-2026-reward}。
% 这些结果表明，显式 rationale 有助于模型在复杂评价标准下组织文本证据和评分判断。
% 然而，相关实验通常将 rationale 生成与提示格式、训练目标、优化策略和推理协议共同改变，因此难以判断性能提升究竟来自 rationale 
% 本身还是完整评价配置。
% 例如，G-Eval 的提升同时依赖 evaluation-step prompting 和特定评分机制；SciRM 的提升也同时包含领域偏好优化、
% reasoning/self-verification 训练和 rationale-first 输出，使 rationale supervision 与 rationale-first inference 
% 的独立作用仍不清楚~\citep{liu-etal-2023-g,sahinuc-etal-2026-reward}。

Existing evaluation methods have introduced rationales through prompting,
supervised distillation and reinforcement learning,
and have generally reported higher human correlation, judge agreement or generalization ability
~\citep{liu-etal-2023-g,kim-etal-2024-prometheus-iclr,sahinuc-etal-2026-reward}.
These results suggest that explicit rationales can help models organize textual evidence
and scoring decisions under complex evaluation criteria.
However, existing experiments often change rationale generation together with prompt format,
training objective, optimization strategy and inference protocol,
making it difficult to determine whether the gains come from rationales themselves
or from the full evaluation configuration.
For example, the gains of G-Eval depend jointly on evaluation-step prompting and its scoring mechanism;
the improvements of SciRM also combine domain-preference optimization,
reasoning/self-verification training and rationale-first outputs,
leaving the independent effects of rationale supervision and rationale-first inference unclear
~\citep{liu-etal-2023-g,sahinuc-etal-2026-reward}.

%=============第三段的思考
% 先用一句话说科学写作评测领域也引入 rationale。
% 然后说有文献发现引入 rationale 没效果
% 甚至是副作用，引出要分清楚两者的独立作用，再然后说有文献专门研究了两者，例如 nlurc 论文就研究了交互作用。
% 然后再说这些研究的问题，有的不是科学写作领域，有的研究的不透彻。

% 为帮助模型依据明确标准组织证据并完成评分，
% 一些科学写作评价器已在评分过程中引入 rationale
% ~\citep{sadallah-etal-2025-good,sahinuc-etal-2026-reward}。
% 然而，rationale 并不必然提升模型表现：rationale augmentation 可能降低任务性能
% ~\citep{zhu-etal-2025-rationales}，而 CoT 的收益主要集中在数学或符号推理任务，
% 在其他任务上并不稳定~\citep{sprague-etal-2025-cot}。
% 为判断性能变化究竟来自训练时学习 rationale，还是推理时先生成 rationale 再评分，
% 现有归因研究考察了这两类设置，但主要面向通用推理或 NLU 任务
% ~\citep{wadhwa-etal-2024-investigating,shi-etal-2025-investigating}。
% 科学写作评价不同于这些任务：模型需要结合领域知识，
% 依据多维 rubric 从文本证据中作出序数判断。
% 因此，现有证据仍无法说明 rationale supervision 与 rationale-first inference
% 在科学写作评价中各自发挥什么作用，以及二者结合时如何相互影响。

To help models organize evidence and assign scores according to explicit criteria,
some scientific-writing evaluators have incorporated rationales into the scoring process
~\citep{sadallah-etal-2025-good,sahinuc-etal-2026-reward}.
However, rationales do not necessarily improve model performance:
rationale augmentation can degrade task performance
~\citep{zhu-etal-2025-rationales},
while CoT gains are concentrated mainly in mathematical or symbolic reasoning tasks
and remain inconsistent elsewhere
~\citep{sprague-etal-2025-cot}.
To determine whether performance changes arise from learning rationales during training
or generating rationales before scoring at inference time,
existing attribution studies have examined both settings,
but primarily on general reasoning or NLU tasks
~\citep{wadhwa-etal-2024-investigating,shi-etal-2025-investigating}.
Scientific-writing evaluation differs from these tasks:
models must integrate domain knowledge and make ordinal judgments
from textual evidence under multidimensional rubrics.
Existing evidence therefore does not establish how rationale supervision
and rationale-first inference individually affect scientific-writing evaluation,
or how they interact when combined.

%=============第四段的思考
%先概括做的事情，包括分析两种训练和推理的接口交互，并对结果做分析，然后提出了
%三个方法。然后下一段写贡献，就是具体描述做了的事情；最后一段写发现，就是
%根据实验结果写发现，把具体的数值结果转成文字描述。

% 本文系统研究 rationale supervision 与 rationale-first inference
% 如何影响基于大模型的科学写作评价。
% 我们构建统一的受控实验框架，在七项科学写作评价任务上
% 交叉比较两种监督目标与两种推理协议。
% 在此基础上，我们从整体性能、样本级预测变化和句子级错误三个层面
% 分析 rationale 的作用，并针对发现的问题设计相应的训练方法。

\paragraph{Our work.}
In this work, we systematically investigate how rationale supervision
and rationale-first inference affect LLM-based scientific-writing evaluation.
We establish a unified controlled framework that crosses two supervision targets
with two inference protocols across seven scientific-writing evaluation tasks.
We then analyze the role of rationales through aggregate performance,
sample-level prediction changes, and sentence-level errors,
and develop targeted training methods based on the identified problems.


% 本文的贡献主要包括三个方面。
% \textbf{(C1)} 我们构建 $2\times2$ 析因设计，将
% score-only supervision（SO）和 rationale-and-score supervision（RS）
% 分别与 direct-score inference（DS）和 rationale-first inference（RF）组合。
% 该设计包含 28 个实验条件，每个条件使用三个随机种子，
% 用于区分监督目标、推理协议及其交互作用。
% \textbf{(C2)} 我们建立多层次分析流程，包括 teacher-forced rationale NLL、
% rationale 质量与一致性评价、配对 DS--RF 预测转移、
% 句子级错误标注以及固定 rationale 后的受控重评分。
% \textbf{(C3)} 我们提出三种针对性方法：
% Regenerated Rationales with Gold-Score Restoration（ReaSon）
% 使用学生模型生成的 rationale 替换教师 rationale 并恢复 gold score；
% Paired-Interface Reweighting（PaIR）通过配对 DS 和 RF 训练视图
% 及等权分支损失缓解训练--推理接口失配；
% Regenerated-Rationale Paired-Interface Reweighting（RePaIR）
% 则将上述两种设计相结合。

\paragraph{Contributions and Findings}
Our contributions are threefold.
\textbf{(C1)} We construct a $2\times2$ factorial design that crosses
score-only supervision (SO) and rationale-and-score supervision (RS)
with direct-score inference (DS) and rationale-first inference (RF).
The design comprises 28 experimental conditions,
each evaluated with three random seeds,
to separate the effects of supervision targets, inference protocols,
and their interaction.
\textbf{(C2)} We develop a multilevel analysis pipeline combining
teacher-forced rationale NLL, rationale quality and consistency assessment,
paired DS--RF prediction transitions, sentence-level error annotation,
and controlled rescoring with fixed rationales.
\textbf{(C3)} We introduce three targeted methods:
Regenerated Rationales with Gold-Score Restoration (ReaSon)
replaces teacher rationales with student-generated ones while restoring gold scores;
Paired-Interface Reweighting (PaIR) trains on paired DS and RF views
with equally weighted branch losses to address training--inference interface mismatch;
and Regenerated-Rationale Paired-Interface Reweighting (RePaIR)
combines the two designs.

% 实验得到三个主要发现。
% \textbf{(F1)} RS 改善了教师 rationale 拟合、生成质量和
% rationale--score 一致性，但这些优势并未转化为更好的评分性能。
% 在序数任务中，RF 在两种监督目标下均增加了严重错误并降低了 QWK；
% 在二分类任务中，其影响则取决于监督目标。
% \textbf{(F2)} 有害转移中的 rationale 通常支持错误分数，
% 且错误主要来自评分标准误用和分数映射。
% 修正错误 rationale 后，158 个所选有害案例中有 142 个恢复为 gold score，
% 表明 rationale 内容能够显著影响这些案例中的后续评分。
% \textbf{(F3)} ReaSon 改善了 RF 序数评分，但未能保持 DS 性能；
% PaIR 在两种推理接口下均优于使用相同数据的 Mix；
% RePaIR 在保持 DS 性能的同时将 RF 的 QWK 从 0.698 提升至 0.735，
% 并在两个接口下保持 100\% 的严格格式有效率。

Our experiments yield three main findings.
\textbf{(F1)} RS improves teacher-rationale fitting,
generated-rationale quality, and rationale--score consistency,
but these gains do not translate into better scoring performance.
On ordinal tasks, RF increases severe errors and lowers QWK
under both supervision objectives, whereas its effects on binary tasks
depend on the supervision objective.
\textbf{(F2)} Rationales in harmful transitions typically support
the incorrect score, with errors arising primarily
from rubric misapplication and score mapping.
Correcting erroneous rationales restores the gold score
in 142 of 158 selected harmful cases,
showing that rationale content can substantially affect subsequent scoring
in these cases.
\textbf{(F3)} ReaSon improves RF ordinal scoring
but does not preserve DS performance;
PaIR outperforms the matched Mix baseline under both inference protocols;
and RePaIR preserves DS performance while increasing RF QWK
from 0.698 to 0.735 and maintaining 100\% strict-format validity
under both interfaces.